\documentclass[12pt,a4paper]{article}
\usepackage[utf8]{inputenc}
\usepackage[T1]{fontenc}
\usepackage[polish]{babel}
\usepackage{amsmath}
\usepackage{graphicx}
\usepackage{booktabs}
\usepackage{geometry}
\usepackage{float}
\usepackage{hyperref}

\geometry{margin=2.5cm}

\title{Sprawozdanie – Projekt 3: Podstawy Gymnasium}
\author{Paweł Adamczyk i Dominika Bujnarowska}
\date{}

\begin{document}

\maketitle

\part{Środowisko dyskretne: CliffWalking}
\textbf{Algorytm:} Q-Learning z tabelą wartości

\section{Opis środowiska}
Wybranym środowiskiem jest \texttt{CliffWalking-v0} z biblioteki Gymnasium. Jest to klasyczny problem planowania na siatce $4 \times 12$ pól, w którym agent porusza się w czterech kierunkach (góra, dół, lewo, prawo), próbując dotrzeć z lewego dolnego rogu siatki do prawego dolnego rogu, omijając krawędź urwiska.

\begin{figure}[H]
    \centering
    \includegraphics[width=0.5\linewidth]{plansza.png}
    \caption{Środowisko Cliff Walking}
    \label{fig:cliff_env}
\end{figure}

\begin{table}[H]
    \centering
    \begin{tabular}{ll}
        \toprule
        \textbf{Parametr} & \textbf{Wartość} \\
        \midrule
        Przestrzeń stanów & Discrete(48) – 48 pól siatki $4 \times 12$ \\
        Przestrzeń akcji & Discrete(4) – góra, dół, lewo, prawo \\
        Nagroda za krok & -1 pkt (każdy bezpieczny krok) \\
        Nagroda za upadek & -100 pkt (teleportacja na start) \\
        Wynik optymalny & -13 pkt (13 kroków bez upadku) \\
        \bottomrule
    \end{tabular}
\end{table}

Środowisko zostało wybrane ze względu na w pełni dyskretny charakter przestrzeni stanów i akcji, co pozwala na bezpośrednie zastosowanie klasycznego algorytmu Q-Learning z tabelą wartości – bez konieczności dyskretyzacji.

\section{Algorytm: Q-Learning}
Q-Learning jest algorytmem uczenia ze wzmocnieniem typu \textit{off-policy}, który uczy się optymalnej funkcji wartości $Q(s, a)$ — przypisującej każdej parze (stan, akcja) oczekiwaną sumę zdyskontowanych nagród.

Reguła aktualizacji ma postać równania Bellmana:
\begin{equation*}
    Q(s, a) \leftarrow Q(s, a) + \alpha \left[ r + \gamma \max_{a'} Q(s', a') \cdot (1 - \text{done}) - Q(s, a) \right]
\end{equation*}

Gdzie:
\begin{itemize}
    \item $\alpha = 0.1$ – współczynnik uczenia (\textit{learning rate})
    \item $\gamma = 0.9$ – współczynnik dyskontowy (wymagany w projekcie)
    \item $r$ – natychmiastowa nagroda ze środowiska
    \item $\max_{a'} Q(s', a')$ – najlepsza znana wartość kolejnego stanu (\textit{greedy})
\end{itemize}

Tablica Q ma rozmiar $48 \times 4$ (stany $\times$ akcje) i jest inicjalizowana zerami. Agent korzysta ze strategii $\varepsilon$-zachłannej ($\varepsilon$-greedy): z prawdopodobieństwem $\varepsilon$ wykonuje losowaną akcję (eksploracja), a z prawdopodobieństwem $1-\varepsilon$ wybiera akcję o najwyższej wartości Q (eksploatacja).

\section{Implementacja}
Kod składa się z dwóch plików:
\begin{itemize}
    \item \texttt{q\_learning\_agent.py} – klasa \texttt{QLearningAgent} zawierająca tabelę Q oraz logikę wyboru akcji i aktualizacji
    \item \texttt{cliff\_walking.ipynb} – notatnik z pętlą trenowania i wizualizacją wyników
\end{itemize}

\subsection*{Główna pętla uczenia}
Funkcja \texttt{train\_agent()} enkapsuluje całą pętlę RL. W każdym epizodzie:
\begin{enumerate}
    \item agent obserwuje stan $s$ i wybiera akcję $a$ (strategia $\varepsilon$-greedy)
    \item środowisko zwraca nagrodę $r$ oraz kolejny stan $s'$
    \item tablica Q jest aktualizowana równaniem Bellmana
    \item po epizodzie epsilon jest mnożony przez współczynnik wygaszania (\textit{decay})
\end{enumerate}
Wartość terminacji przekazywana do równania aktualizacji rozróżnia stan końcowy (wytrenowany, bez przyszłych nagród) od przerwania przez limit kroków (\textit{truncated}), co zapewnia poprawną propagację wartości Q.

\section{Eksperymenty}
Przeprowadzono eksperyment porównujący trzy strategie wygaszania parametru $\varepsilon$ na przestrzeni 500 epizodów. Wszystkie pozostałe hiperparametry ($\alpha=0.1$, $\gamma=0.9$) pozostawały stałe. Na wykresach widoczna jest średnia krocząca z okna 20 epizodów.

\begin{table}[H]
    \centering
    \begin{tabular}{llcl}
        \toprule
        \textbf{Konfiguracja} & \textbf{start $\varepsilon$} & \textbf{decay} & \textbf{Wynik (ostatnie 50 ep.)} \\
        \midrule
        Stały $\varepsilon = 0.1$ & 0.1 & 1.0 (brak) & $\sim -40$ do $-60$ (duża wariancja) \\
        Szybkie wygaszanie & 1.0 & 0.95 & $\sim -13$ (stabilnie optymalna) \\
        Wolne wygaszanie & 1.0 & 0.99 & $\sim -13$ do $-20$ (wolniejsza konwergencja) \\
        \bottomrule
    \end{tabular}
\end{table}

\begin{figure}[H]
    \centering
    \includegraphics[width=0.95\textwidth]{wykres1.png}
    \caption{Porównanie krzywych uczenia i przebiegu wygaszania $\varepsilon$ dla trzech strategii eksploracji w CliffWalking.}
\end{figure}

\section{Opis procesu uczenia}
\subsection*{Faza eksploracji (epizody 0–100)}
Na początku procesu uczenia agent działa niemal losowo – wysoka wartość $\varepsilon$ powoduje częste spadki z urwiska i nagrody rzędu $-100$. W tym czasie tablica Q jest intensywnie zapełniana wstępnymi oszacowaniami. Krzywa uczenia wykazuje dużą wariancję.

\subsection*{Faza przejściowa (epizody 100–300)}
W miarę wygaszania $\varepsilon$ agent coraz częściej korzysta z wyuczonej tablicy Q. Wyniki stabilizują się i wyraźnie rosną. Szybkie wygaszanie (\texttt{decay=0.95}) wyprzedza wolniejsze – agent wcześniej przechodzi do eksploatacji poznanej ścieżki.

\subsection*{Faza eksploatacji (epizody 300–500)}
Po osiągnięciu $\varepsilon \approx 0$ obie wersje z wygaszaniem osiągają wynik bliski optimum ($-13$ pkt). Strategia ze stałym $\varepsilon=0.1$ nigdy nie osiąga pełnej stabilności ze względu na 10\% losowości.

\section{Wnioski}
\begin{itemize}
    \item Q-Learning skutecznie rozwiązał problem CliffWalking, osiągając nagrodę $-13$.
    \item Strategia wygaszania $\varepsilon$ ma krytyczny wpływ na jakość uczenia.
    \item Współczynnik $\gamma=0.9$ pozwolił na poprawne zaplanowanie ścieżki omijającej urwisko.
\end{itemize}

\newpage
\part{Środowisko ciągłe: CartPole-v1}

\section{Opis środowiska i cel}

Środowisko \texttt{CartPole-v1} z biblioteki Gymnasium modeluje klasyczny problem
sterowania odwróconym wahadłem. Na wózku poruszającym się wzdłuż poziomego toru
osadzony jest kij, który agent musi utrzymać w pionie — jedyną dostępną akcją
jest pchnięcie wózka w lewo lub w prawo.

\begin{table}[h]
    \centering
    \begin{tabular}{ll}
        \toprule
        \textbf{Parametr} & \textbf{Wartość} \\
        \midrule
        Przestrzeń obserwacji & Box(4,) – 4 zmienne rzeczywiste (ciągła) \\
        Przestrzeń akcji & Discrete(2) – pchnięcie w lewo lub w prawo \\
        Nagroda za krok & +1 pkt (każdy krok utrzymania kija) \\
        Koniec epizodu & kąt kija $> 12°$ lub pozycja wózka $> 2{,}4$ \\
        Wynik maksymalny & 500 pkt \\
        \bottomrule
    \end{tabular}
    \caption{Parametry środowiska CartPole-v1.}
\end{table}

\begin{figure}[h]
    \centering
    \includegraphics[width=0.75\linewidth]{cartpole-env.png}
    \caption{Środowisko CartPole-v1.}
\end{figure}

Kluczowa różnica względem \texttt{CliffWalking}: przestrzeń obserwacji jest
\textbf{ciągła}. Wektor stanu zawiera cztery liczby rzeczywiste:
pozycję wózka, jego prędkość, kąt kija i prędkość kątową kija.
Klasyczny Q-Learning z tablicą wymaga skończonej przestrzeni stanów —
konieczna jest zatem \textit{dyskretyzacja} lub zastosowanie głębokiej sieci
neuronowej.

W tej części sprawozdania prezentujemy oba podejścia:
\begin{enumerate}
    \item \textbf{Q-Learning z dyskretyzacją} — punkt odniesienia; pokazuje
    ograniczenia podejścia tablicowego dla ciągłych przestrzeni obserwacji.
    \item \textbf{Deep Q-Network (DQN)} — sieć neuronowa aproksymuje funkcję
    $Q(s,a)$ bezpośrednio na ciągłym wektorze stanu.
\end{enumerate}

% ============================================================
\section{Algorytm 1: Q-Learning z dyskretyzacją}
% ============================================================

\subsection{Idea dyskretyzacji}

Każda z czterech zmiennych obserwacji jest mapowana na dyskretne kubełki
(\textit{bins}) przy pomocy funkcji \texttt{np.digitize}. Przyjęto 6 progów
na wymiar, co tworzy 7 kubełków i daje $7^4 = 2401$ możliwych stanów
dyskretnych.

\begin{table}[h]
    \centering
    \begin{tabular}{lll}
        \toprule
        \textbf{Zmienna} & \textbf{Zakres} & \textbf{Liczba progów} \\
        \midrule
        Pozycja wózka & $[-4{,}8;\; 4{,}8]$ & 6 \\
        Prędkość wózka & $[-4;\; 4]$ & 6 \\
        Kąt kija (rad) & $[-0{,}418;\; 0{,}418]$ & 6 \\
        Prędkość kątowa kija & $[-4;\; 4]$ & 6 \\
        \bottomrule
    \end{tabular}
    \caption{Parametry dyskretyzacji przestrzeni obserwacji.}
\end{table}

Tablica Q ma rozmiar $2401 \times 2$ i jest inicjalizowana zerami.
Agent korzysta ze strategii $\varepsilon$-zachłannej identycznej jak
w części pierwszej sprawozdania (CliffWalking), z hiperparametrami:
$\alpha = 0{,}1$, $\varepsilon_0 = 1{,}0$, \textit{decay} $= 0{,}995$,
$\varepsilon_{\min} = 0{,}01$.

\subsection{Eksperyment: wpływ współczynnika dyskontowego $\gamma$}

Przeprowadzono trening dla trzech wartości $\gamma \in \{0{,}5;\; 0{,}9;\; 0{,}99\}$
przez 1500 epizodów. Metrykę porównawalności stanowi całkowita zdyskontowana
nagroda z pierwszych 1000 epizodów przy $\gamma = 0{,}9$:

\begin{equation*}
    G = \sum_{i=0}^{999} \gamma^i \cdot R_i, \quad \gamma = 0{,}9
\end{equation*}

\begin{table}[h]
    \centering
    \begin{tabular}{p{2.5cm}p{3.5cm}p{3.5cm}p{4cm}}
        \toprule
        \textbf{$\gamma$} & \textbf{Wynik (ost. 100 ep.)} & \textbf{Zdysk. nagroda $G$} & \textbf{Charakterystyka} \\
        \midrule
        $0{,}5$  & $\sim 80{-}130$ pkt  & 209 & duża wariancja, niestabilna konwergencja \\
        $0{,}9$  & $\sim 90{-}120$ pkt  & \textbf{229} & najgładszy wzrost, najwyższe $G$ \\
        $0{,}99$ & $\sim 100{-}140$ pkt & 182 & wysokie szczyty, niestabilna propagacja \\
        \bottomrule
    \end{tabular}
    \caption{Wyniki Q-Learningu z dyskretyzacją dla trzech wartości $\gamma$.
    Zdyskontowana nagroda $G$ liczona z pierwszych 1000 epizodów przy $\gamma=0{,}9$.}
\end{table}

\begin{figure}[h]
    \centering
    \includegraphics[width=0.95\textwidth]{cartpole-q.png}
    \caption{Krzywe uczenia Q-Learningu z dyskretyzacją (średnia krocząca, okno 50 epizodów)
    dla trzech wartości $\gamma$. Linia przerywana oznacza maksymalny wynik 500~pkt.}
    \label{fig:ql_gamma}
\end{figure}

\subsection{Opis procesu uczenia i interpretacja $\gamma$}

W fazie eksploracji (epizody 0–400) wszystkie konfiguracje zachowują się
niemal identycznie — wysoki $\varepsilon$ powoduje głównie losowe akcje
i nagrody rzędu 20–40 pkt. Różnice ujawniają się po epizodzie 400,
gdy $\varepsilon$ opada poniżej 0{,}1.

\begin{itemize}
    \item \textbf{$\gamma = 0{,}5$ (myślenie krótkoterminowe):}
    Agent silnie premiuje nagrody z najbliższego kroku.
    Reaguje szybko na chwilowe zagrożenia, lecz nie buduje strategii
    długofalowej. Skutkuje to dużą wariancją i brakiem stabilnego
    optimum — mimo że lokalne reakcje są poprawne, agent nie potrafi
    zaplanować utrzymania kija przez dłuższy czas.
    Zdyskontowana nagroda $G = 209$ jest najniższa spośród
    konfiguracji z wygaszaniem.

    \item \textbf{$\gamma = 0{,}9$ (wymagany projektem):}
    Balans między nagrodami natychmiastowymi a odległymi.
    Daje najgładszą krzywą uczenia i najwyższą zdyskontowaną nagrodę
    $G = 229$. W połączeniu z dyskretyzacją jest to najstabilniejsza
    konfiguracja — agent konsekwentnie poprawia wyuczoną politykę
    bez wyraźnych regresji.

    \item \textbf{$\gamma = 0{,}99$ (planowanie długoterminowe):}
    Teoretycznie powinno sprzyjać utrzymaniu kija przez 500 kroków,
    jednak przy niedokładnej dyskretyzacji błędy w szacowaniu
    odległych stanów kumulują się podczas propagacji wstecznej.
    Aktualizacje Bellmana stają się niestabilne — stąd pomimo
    chwilowych szczytów powyżej 140 pkt, zdyskontowana nagroda
    $G = 182$ jest \emph{najniższa} ze wszystkich trzech wartości.
\end{itemize}

\subsection{Ograniczenia dyskretyzacji}

Po epizodzie 900 wszystkie konfiguracje osiągają plateau w przedziale
100–140 pkt, niezależnie od wartości $\gamma$. Wynik ten nie pochodzi
z błędu implementacji — jest \textbf{sufitem dyskretyzacji}:
przy 7 kubełkach na wymiar zbyt wiele różnych stanów fizycznych trafia
do tego samego indeksu tablicy Q, uniemożliwiając dalsze rozróżnianie
akcji. Zwiększenie liczby kubełków poprawia dokładność, ale
wykładniczo zwiększa liczbę stanów do odwiedzenia ($N^4$ dla $N$ kubełków),
wydłużając czas eksploracji ponad praktyczne możliwości.

Rozwiązaniem jest bezpośrednia aproksymacja $Q(s,a)$ siecią neuronową,
co eliminuje konieczność dyskretyzacji.

% ============================================================
\section{Algorytm 2: Deep Q-Network (DQN)}
% ============================================================

\subsection{Idea i kluczowe mechanizmy}

DQN zastępuje tablicę Q siecią neuronową, która dla dowolnego ciągłego
wektora stanu $s$ zwraca wartości $Q(s,a)$ dla każdej akcji.
Implementacja opiera się na oficjalnym tutorialu PyTorch i zawiera
dwa mechanizmy stabilizujące uczenie:

\begin{enumerate}
    \item \textbf{Experience Replay:} przejścia $(s, a, r, s')$ są
    zapisywane do bufora (\textit{Replay Memory}) o pojemności 10\,000
    i losowane w mini-batchach po 256 próbek. Przerywa to korelację
    czasową między kolejnymi aktualizacjami.

    \item \textbf{Target Network z soft update:} zamiast kopiowania
    wag co $N$ kroków (hard update), target network jest aktualizowana
    co krok według reguły:
    \begin{equation*}
        \theta_{\text{target}} \leftarrow \tau \cdot \theta_{\text{policy}}
        + (1-\tau) \cdot \theta_{\text{target}}, \quad \tau = 0{,}01
    \end{equation*}
    Wagi target network zmieniają się \emph{płynnie}, co eliminuje
    skokowe destabilizacje charakterystyczne dla hard update.
\end{enumerate}

Jako funkcję straty użyto Huber loss (\texttt{SmoothL1Loss}),
a gradienty są przycinane do wartości bezwzględnej 100
(\textit{gradient clipping}).

\subsection{Dwie architektury sieci}

Porównano dwie konfiguracje architektury sieci:

\begin{table}[h]
    \centering
    \begin{tabular}{lll}
        \toprule
        \textbf{Architektura} & \textbf{Warstwy} & \textbf{Parametry} \\
        \midrule
        \texttt{DQN\_128x128} & $4 \to 128 \to 128 \to 2$ & 18\,818 \\
        \texttt{DQN\_128x64}  & $4 \to 128 \to 64  \to 2$ & 9\,474  \\
        \bottomrule
    \end{tabular}
    \caption{Porównane architektury sieci neuronowej. Aktywacja ReLU po każdej
    warstwie ukrytej. Optymalizator AdamW z \texttt{amsgrad=True}.}
\end{table}

\subsection{Przeszukiwanie hiperparametrów}

Przetestowano 3 konfiguracje hiperparametrów dla każdej architektury
(łącznie 6 treningów po 600 epizodów). Zmieniano jeden parametr
względem konfiguracji bazowej, aby ocenić jego izolowany wpływ.

\begin{table}[h]
    \centering
    \begin{tabular}{llccc}
        \toprule
        \textbf{Konfiguracja} & \textbf{Zmiana} & \textbf{BATCH} & \textbf{LR} & \textbf{EPS\_DECAY} \\
        \midrule
        bazowa & — & 256 & $5 \cdot 10^{-4}$ & 500 \\
        wolniejszy decay & dłuższa eksploracja & 256 & $5 \cdot 10^{-4}$ & 1000 \\
        niższe LR & ostrożniejsze aktualizacje & 256 & $1 \cdot 10^{-4}$ & 500 \\
        \bottomrule
    \end{tabular}
    \caption{Siatka hiperparametrów. Pozostałe parametry stałe:
    $\gamma = 0{,}99$, $\varepsilon_0 = 0{,}9$, $\varepsilon_{\min} = 0{,}05$,
    TAU $= 0{,}01$.}
\end{table}

\textbf{Metryka optymalizacji:} całkowita zdyskontowana nagroda
z 600 epizodów przy $\gamma = 0{,}99$:
\begin{equation*}
    G_{\text{DQN}} = \sum_{i=0}^{599} \gamma^i \cdot D_i, \quad \gamma = 0{,}99
\end{equation*}
gdzie $D_i$ to czas trwania $i$-tego epizodu (równy sumie nagród, bo nagroda = +1 co krok).

\subsection{Wyniki przeszukiwania hiperparametrów}

\begin{table}[h]
    \centering
    \begin{tabular}{llr}
        \toprule
        \textbf{Architektura} & \textbf{Konfiguracja} & \textbf{Zdysk. nagroda $G_{\text{DQN}}$} \\
        \midrule
        \texttt{DQN\_128x128} & bazowa           & \textbf{17\,902} \\
        \texttt{DQN\_128x64}  & bazowa           & 16\,747 \\
        \texttt{DQN\_128x128} & wolniejszy decay & 16\,317 \\
        \texttt{DQN\_128x64}  & wolniejszy decay & 10\,725 \\
        \texttt{DQN\_128x128} & niższe LR        & 2\,046  \\
        \texttt{DQN\_128x64}  & niższe LR        & 1\,221  \\
        \bottomrule
    \end{tabular}
    \caption{Ranking konfiguracji DQN według zdyskontowanej nagrody $G_{\text{DQN}}$
    z 600 epizodów ($\gamma = 0{,}99$). Najlepsza konfiguracja: \texttt{DQN\_128x128} bazowa.}
\end{table}

\begin{figure}[H]
    \centering
    \includegraphics[width=0.95\textwidth]{dqn-128x64-hp.png}
    \caption{Krzywe uczenia dla architektury \texttt{DQN\_128x64} — porównanie
    trzech konfiguracji hiperparametrów.}
    \label{fig:dqn_128x64}
\end{figure}

\subsection{Interpretacja wyników}

Wyniki wykazują wyraźny wzorzec:

\begin{itemize}
    \item \textbf{Konfiguracja bazowa ($G \approx 17\,900 / 16\,700$):}
    Domyślne parametry z tutoriala PyTorch okazały się najlepsze.
    $\varepsilon$ opada do minimum po $\approx 1500$ krokach
    (kilkanaście epizodów), co pozwala agentowi szybko przejść
    do eksploatacji wyuczonej polityki.

    \item \textbf{Wolniejszy decay ($G \approx 16\,300 / 10\,700$):}
    Wydłużenie eksploracji do $\approx 3000$ kroków (EPS\_DECAY=1000)
    pozwala agentowi dokładniej poznać przestrzeń stanów, ale opóźnia
    konwergencję. Dla \texttt{DQN\_128x64} różnica jest znaczna
    ($-36\%$), co sugeruje że mniejsza sieć potrzebuje mniej danych
    do nauki stabilnej polityki.

    \item \textbf{Niższe LR ($G \approx 2\,046 / 1\,221$):}
    Obniżenie LR z $5 \cdot 10^{-4}$ do $1 \cdot 10^{-4}$ drastycznie
    pogarsza wyniki — przy 600 epizodach sieć nie zdąża
    nauczyć się użytecznej polityki. Mechanizm soft update
    z $\tau = 0{,}01$ wymaga stosunkowo wysokiego LR, aby target network
    nadążał za poprawiającą się siecią główną.
\end{itemize}

% ============================================================
\section{Porównanie algorytmów: Q-Learning vs DQN}
% ============================================================

\begin{figure}[H]
    \centering
    \includegraphics[width=0.95\textwidth]{cartpole-final-comparison.png}
    \caption{Porównanie końcowe: Q-Learning z dyskretyzacją ($\gamma = 0{,}9$)
    vs najlepsza konfiguracja DQN (\texttt{DQN\_128x128}, bazowa, $\gamma = 0{,}99$).
    Średnia krocząca z okna 50 epizodów.}
    \label{fig:final_comparison}
\end{figure}

\begin{table}[H]
    \centering
    \begin{tabular}{p{4.5cm}p{3cm}p{3cm}p{2.5cm}}
        \toprule
        \textbf{Algorytm} & \textbf{Repr.\ stanów} & \textbf{Wynik (ost.\ 100 ep.)} & \textbf{$G$ (1000 ep.)} \\
        \midrule
        Q-Learning, $\gamma = 0{,}5$ & dyskretna & $80{-}130$ pkt  & 209 \\
        Q-Learning, $\gamma = 0{,}9$ & dyskretna & $90{-}120$ pkt  & \textbf{229} \\
        Q-Learning, $\gamma = 0{,}99$ & dyskretna & $100{-}140$ pkt & 182 \\
        \midrule
        DQN 128x128 [bazowa]          & ciągła & $\approx 500$ pkt & \multirow{2}{*}{$\sim 17\,900$} \\
        \quad $\gamma = 0{,}99$       &        &                   & \\
        \bottomrule
    \end{tabular}
    \caption{Zbiorcze porównanie algorytmów. $G$ dla Q-Learningu liczone przy $\gamma=0{,}9$
    z 1000 epizodów; dla DQN przy $\gamma=0{,}99$ z 600 epizodów.}
\end{table}

Różnica między algorytmami jest fundamentalna:

\begin{itemize}
    \item Q-Learning z dyskretyzacją utyka na plateau $\sim 140$ pkt
    (28\% maksimum) niezależnie od $\gamma$ i liczby epizodów.
    Jest to ograniczenie reprezentacji, nie algorytmu.

    \item DQN osiąga wynik $\approx 500$ pkt (maksimum środowiska)
    po kilkuset epizodach. Sieć neuronowa aproksymuje $Q(s,a)$
    na ciągłym wektorze stanu bez utraty informacji.

    \item Kluczowym ulepszeniem DQN względem naiwnego Q-Learningu
    z siecią jest \textbf{soft update} target network ($\tau = 0{,}01$
    co krok). W porównaniu z hard update co $N$ kroków, wagi
    target network zmieniają się płynnie — eliminuje to skokowe
    destabilizacje, które prowadziły do \textit{catastrophic forgetting}
    (gwałtownego spadku wyniku z $\sim 450$ do $\sim 60$ pkt
    obserwowanego przy hard update).
\end{itemize}

% ============================================================
\section{Wizualizacja wytrenowanego agenta}
% ============================================================

Po zakończeniu treningu najlepsza sieć (\texttt{DQN\_128x128}, konfiguracja bazowa)
jest zapisywana do pliku \texttt{best\_policy\_net.pt}.
Plik \texttt{cartpole\_visualize.py} ładuje zapisany model i uruchamia
animację w czasie rzeczywistym przy użyciu biblioteki \texttt{pygame}.

Agent działa czysto zachłannie — bez eksploracji ($\varepsilon = 0$).
Wyświetlane są: aktualny kąt kija, pozycja wózka, numer kroku
oraz suma nagród. Sterowanie: klawisz \textbf{R} resetuje epizod,
\textbf{Q} kończy wizualizację.

% ============================================================
\section{Wnioski}
% ============================================================

\begin{enumerate}
    \item \textbf{Dyskretyzacja tworzy sufit.}
    Q-Learning z tablicą osiąga $\sim 100{-}140$ pkt (20–28\% maksimum)
    niezależnie od $\gamma$. Sufit wynika z utraty informacji przy
    mapowaniu stanów na kubełki — nie jest błędem implementacji.

    \item \textbf{Wpływ $\gamma$ na Q-Learning.}
    $\gamma = 0{,}9$ daje najwyższą zdyskontowaną nagrodę ($G = 229$)
    i najstabilniejsze uczenie. $\gamma = 0{,}99$ jest
    kontr-produktywne przy dyskretyzacji: błędy kubełkowania kumulują
    się przy propagacji długoterminowych wartości Q, obniżając $G$
    do 182. $\gamma = 0{,}5$ daje dużą wariancję i pośrednią wartość
    $G = 209$.

    \item \textbf{DQN eliminuje ograniczenia dyskretyzacji.}
    Najlepsza konfiguracja (\texttt{DQN\_128x128}, bazowa) osiąga
    zdyskontowaną nagrodę $G_{\text{DQN}} = 17\,902$ — o rzędy
    wielkości wyższą niż Q-Learning — i regularnie osiąga maksymalny
    wynik 500 pkt.

    \item \textbf{Soft update jest kluczowy dla stabilności.}
    $\tau = 0{,}01$ co krok eliminuje catastrophic forgetting.
    Przy hard update obserwowano gwałtowne restarty uczenia;
    przy soft update krzywa rośnie monotonicznie.

    \item \textbf{LR ma dominujący wpływ na wyniki DQN.}
    Obniżenie LR z $5 \cdot 10^{-4}$ do $1 \cdot 10^{-4}$
    redukuje $G$ o $\sim 89\%$ dla obu architektur —
    jest to najważniejszy hiperparametr przy 600 epizodach treningu.
    Prędkość eksploracji (EPS\_DECAY) ma umiarkowany wpływ.

    \item \textbf{Architektura sieci ma drugorzędne znaczenie dla CartPole.}
    \texttt{DQN\_128x128} ($G = 17\,902$) nieznacznie wyprzedza
    \texttt{DQN\_128x64} ($G = 16\,747$) przy konfiguracji bazowej.
    Różnica wynosi $\sim 7\%$ — CartPole jest wystarczająco proste,
    by obie architektury skutecznie je rozwiązywały.
\end{enumerate}

\end{document}